\chapter{Methodology}
\label{ch:methodology}

\section{System overview}
\label{sec:overview}

The system is a five-layer stack (Figure~\ref{fig:architecture}) built around one
idea: every claim an agent makes should be traceable to something it was actually
given. A document goes in, an extraction agent turns it into a validated scenario
object, four agents interact under an explicit state machine, and every step is
written to a JSON log that later analysis reads without re-invoking the model. The
Python backend is 33 modules and roughly 4,900 lines, with a 92-case test suite;
the API and frontend add a further 374 and 3,800 lines respectively.

\begin{figure}[htbp]
\centering
\begin{tikzpicture}[font=\footnotesize,
    layer/.style={draw=f21ink,line width=0.7pt,rounded corners=2pt,
                  minimum width=12.4cm,minimum height=1.02cm,align=center,fill=f21mist},
    gap/.style={draw=f21grey,dashed,line width=0.7pt,rounded corners=2pt,
                minimum width=12.4cm,minimum height=0.78cm,align=center,fill=white},
    tag/.style={anchor=east,f21grey,font=\scriptsize\itshape}]

  \node[layer] (l1) at (0,0)
    {\textbf{L1 Presentation} --- React/TypeScript client: upload, live streamed negotiation,\\
     case browser, analytics, risk screen};
  \node[layer,below=3.5mm of l1] (l2)
    {\textbf{L2.5 API} --- FastAPI: extraction endpoint, server-sent-event negotiation stream,\\
     risk assessment, GPU-tunnel status and on-demand reconnect};
  \node[layer,below=3.5mm of l2] (l3)
    {\textbf{L2 Orchestration} --- LangGraph \texttt{StateGraph}: deterministic turn order,\\
     conditional routing, per-super-step streaming};
  \node[layer,below=3.5mm of l3] (l4)
    {\textbf{L3 Agents} --- Contracting Authority, Aggrieved Bidder, Court, Summary, Extraction\\
     (Llama 3.1 8B via Ollama; local CPU or tunnelled A10G GPU)};
  \node[layer,below=3.5mm of l4] (l5)
    {\textbf{L5 Persistence} --- structured JSON logs per run; batch summaries with\\
     provenance and compliance counters};
  \node[gap,below=3.5mm of l5] (lg)
    {\textbf{L4 Knowledge / RAG} --- vector retrieval over legislation and case law\\
     \emph{--- deliberately not built; see \S\ref{sec:alternatives}}};

  \node[tag] at ($(l1.west)+(-0.15,0)$) {browser};
  \node[tag] at ($(l3.west)+(-0.15,0)$) {process};
  \node[tag] at ($(l4.west)+(-0.15,0)$) {model};
\end{tikzpicture}
\caption[System architecture]{The five-layer architecture. The dashed layer is
reported because it was scoped out, not because it exists: the extraction agent is
a single-pass synthesis step and must not be mistaken for retrieval.}
\label{fig:architecture}
\end{figure}

\section{Grounding the agents in the client's own material}
\label{sec:grounding}

An early version of this system used interest lists I had written myself, which
made the agents plausible and unfalsifiable in equal measure. They were replaced
with the taxonomies in Fusion21's Doc 1, so that each agent's stated interests
trace to a named category in a document the client wrote rather than to my
intuition about what a procurement officer cares about. Doc 1 gives seven
categories for the contracting authority, six for the supplier, and six for the
judiciary; Table~\ref{tab:agent-config} summarises the resulting configuration.

Doc 1 also supplies per-case behavioural modifiers, its ``practical
considerations'' tables, which vary by dispute rather than being fixed properties
of a role. These are modelled as optional \texttt{CAProfile} and
\texttt{BidderProfile} objects attached to a scenario and rendered into the system
prompt at construction. Both default to unset and render no text, which keeps
every result obtained before profiles existed directly comparable. That
decision has a structural consequence: agents cannot be module-level singletons,
because their system prompts now depend on the scenario, so they are built per
dispute and bound into the graph nodes at compile time.

\begin{table}[htbp]
\centering
\footnotesize
\caption[Agent configuration]{Agent configuration. Temperatures were set by role
function rather than tuned: the adjudicator, summariser and extractor are held low
because consistency matters more than variety, the challenger is highest because a
negotiation in which nobody proposes anything new stalls. The two negotiating
agents escalate their own temperature on a repetition retry, to 0.8 and 0.9
respectively.}
\label{tab:agent-config}
\begin{tabular}{@{}llcp{5.4cm}@{}}
\toprule
\textbf{Agent} & \textbf{Interest source} & \textbf{Temp.} & \textbf{Structured output} \\
\midrule
Contracting Authority & Doc 1, 7 CA categories & 0.3 &
  \texttt{PreNegotiationStatement}, \texttt{RoundResponse}, \texttt{WinStatement} \\
Aggrieved Bidder & Doc 1, 6 supplier categories & 0.4 &
  \texttt{PreNegotiationStatement}, \texttt{RoundResponse}, \texttt{WinStatement} \\
Court / Judge & Doc 1, 6 judiciary categories & 0.2 &
  \texttt{ComplianceAssessment} \\
Summary (non-negotiating) & --- & 0.2 &
  \texttt{NegotiationSummary} \\
Extraction & --- & 0.2 &
  \texttt{DisputeScenario} \\
\bottomrule
\end{tabular}
\end{table}

The fourth agent was added at my supervisor's suggestion and does not negotiate.
It reads the completed transcript and produces a plain-English account of what
happened, on the reasoning that a system whose output only a procurement lawyer
can read has not actually reduced anybody's costs. Section~\ref{sec:readability}
reports whether it succeeds, and the answer is not flattering.

\section{The Court agent: conditional verification}
\label{sec:court-design}

The Court agent assesses process compliance only. This mirrors how UK courts
actually review procurement decisions: the question is whether the authority
followed a lawful, rational and procedurally correct process, not whether the
right bidder won. Getting that agent to behave took four revisions, each driven by
a measured failure rather than by taste (Table~\ref{tab:court-versions}).

\begin{table}[htbp]
\centering
\footnotesize
\caption[Court prompt revisions]{Four revisions of the Court agent's prompt. Each
row's failure mode is what motivated the next row.}
\label{tab:court-versions}
\begin{tabular}{@{}llp{7.0cm}@{}}
\toprule
\textbf{Version} & \textbf{Instruction} & \textbf{Observed behaviour} \\
\midrule
V1 & Judgement-based, no explicit check &
  Manifest-error findings tracked whether the CA conceded in dialogue, not the facts \\
V2 & Always verify any calculation &
  Fixed the numeric cases, but invented point systems on qualitative disputes that had none \\
V3 & Conditional: compute if a real formula exists, otherwise qualitative review &
  Restored non-fabricating judgement while keeping numeric accuracy \\
V4 & V3 with the judiciary interest block replaced by Doc 1's six categories &
  Never deadlocks on any case; see \S\ref{sec:ablation} for whether that is good news \\
\bottomrule
\end{tabular}
\end{table}

V4 is derived from V3 by a single verified string substitution of one named block,
and the module raises at import if that anchor ever stops matching. The two
prompts therefore differ in exactly one variable, which is what makes the ablation
in Section~\ref{sec:ablation} a controlled comparison rather than a
before-and-after anecdote.

The decision procedure itself is shown in Figure~\ref{fig:court-flow}. Its
non-obvious element is the middle branch. A scenario that states a formula and a
correction, but not the original inputs, is not computable, and an earlier version
of the prompt allowed the model to close that gap by assuming the missing figures.
Doing so produced arithmetic that was internally valid and evidentially worthless.
Step 1's gate was therefore tightened to require the formula \emph{and} every
input value it needs, with an explicit statement that a stated correction
describes an outcome rather than an input, backed by a ban on writing down an
unstated number even when hedged as illustrative or assumed. The hedge is not a
mitigation; it is the exact linguistic form the failure took.

\begin{figure}[htbp]
\centering
\begin{tikzpicture}[font=\footnotesize,>=Latex,
  box/.style={draw=f21ink,rounded corners=2pt,align=center,
              inner sep=4pt,fill=white,line width=0.7pt},
  dec/.style={draw=f21ink,align=center,inner sep=4pt,fill=f21mist,line width=0.7pt},
  out/.style={draw=f21green,line width=0.9pt,rounded corners=2pt,align=center,
              inner sep=4pt,fill=f21green!8,text width=4.1cm},
  bad/.style={draw=f21amber,line width=0.9pt,rounded corners=2pt,align=center,
              inner sep=4pt,fill=f21amber!8,text width=4.1cm}]

  \node[box]           (in) at ( 0.0, 5.4) {Round transcript + scenario};
  \node[dec,text width=8.4cm] (q1) at ( 0.0, 4.0)
      {\textbf{Step 1.} Does the scenario state a formula\\
       \emph{and} every input value that formula needs?};

  \node[out] (a) at (-3.4, 2.0) {\textbf{Step 2A}\\
      Recompute exactly, show working.\\ Mismatch $\Rightarrow$ manifest error};
  \node[out] (b) at ( 3.4, 2.0) {\textbf{Step 2B}\\
      Qualitative rational-basis review.\\ No invented numbers};
  \node[bad] (c) at ( 0.0,-0.5) {\textbf{Formula stated, inputs missing}\\
      Do not assume, estimate or\\ ``illustrate'' the gap};

  \draw[->] (in) -- (q1);
  \draw[->] (q1.south) -- ++(0,-0.35)
      -| node[pos=0.28,above,font=\scriptsize] {both present} (a.north);
  \draw[->] (q1.south) -- ++(0,-0.35)
      -| node[pos=0.28,above,font=\scriptsize] {neither} (b.north);
  \draw[->] (q1.south) -- node[pos=0.72,right,font=\scriptsize] {formula only} (c.north);
  \draw[->,f21amber,line width=1pt] (c.east)
      -| node[pos=0.22,below,font=\scriptsize] {fall back to 2B} (b.south);
\end{tikzpicture}
\caption[Court agent decision procedure]{The Court agent's gating procedure. The
amber path is the one that had to be added after the Woods v Milton Keynes case
exposed it: a formula plus a stated correction is not a computable scenario, and
treating it as one produced fabricated inputs.}
\label{fig:court-flow}
\end{figure}

The same discipline is applied to legal citation. The agent may cite a provision
only if it appears in the scenario, in a party's statement that round, or in the
grounding material given in its own system prompt. A correct general point with no
citation is explicitly preferred to a specific-sounding invented one, on the basis
that a fabricated regulation number is the same category of error as a fabricated
sub-score. Section~\ref{sec:citation} measures whether the instruction holds.

\section{Orchestration}
\label{sec:orchestration}

Turn order is a LangGraph \texttt{StateGraph} (Figure~\ref{fig:graph}) rather than
a conversational free-for-all. Both parties state their interests, goals and BATNA
before anyone negotiates; then the authority and the bidder alternate, with the
Court assessing compliance after every round. A single conditional edge decides
whether to loop.

\begin{figure}[htbp]
\centering
\begin{tikzpicture}[font=\footnotesize,>=Latex,
  n/.style={draw=f21ink,rounded corners=2pt,minimum height=7.5mm,inner xsep=5pt,
            fill=white,line width=0.7pt},
  c/.style={draw=f21green,rounded corners=2pt,minimum height=7.5mm,inner xsep=5pt,
            fill=f21green!8,line width=0.9pt}]
  \node[n] (pre) at (0,0) {pre\_negotiation};
  \node[n,right=8mm of pre] (ca) {ca\_round};
  \node[n,right=8mm of ca] (bid) {bidder\_round};
  \node[c,right=8mm of bid] (crt) {court\_check};
  \node[n,right=10mm of crt] (win) {win\_statements};
  \node[n,below=7mm of win] (sum) {summary};
  \node[right=8mm of sum] (end) {\textsc{end}};

  \draw[->] (pre) -- (ca);
  \draw[->] (ca) -- (bid);
  \draw[->] (bid) -- (crt);
  \draw[->] (crt) -- node[above,font=\scriptsize,align=center]
      {resolved\\ or round $\geq$ max} (win);
  \draw[->] (win) -- (sum);
  \draw[->] (sum) -- (end);
  \draw[->,f21green] (crt.south) -- ++(0,-0.75) -| node[below,pos=0.25,font=\scriptsize]
      {\textcolor{f21green}{continue}} (ca.south);
\end{tikzpicture}
\caption[Negotiation state graph]{The negotiation graph. Everything hinges on
\texttt{court\_check}: it is the only node that can set \texttt{resolved}, and
therefore the only escape from the loop other than exhausting the round limit.
Section~\ref{sec:baselines} removes it and measures the consequence.}
\label{fig:graph}
\end{figure}

Resolution is decided by the Court's \texttt{recommended\_action}, drawn from a
closed vocabulary: continue negotiation, re-evaluation, no remedy with the
decision standing, or damages. Anything other than ``continue'' resolves the
dispute. Exhausting the round limit without resolution is recorded as deadlock and
escalation to formal proceedings. Constraining this vocabulary is a legal
requirement rather than an implementation convenience --- Doc 3 is explicit that
award-stage disputes cannot be settled by splitting the difference --- and the
constraint is enforced rather than merely documented: the downstream
recommendation module raises if a batch contains any outcome outside the set,
which was verified by deliberately injecting the string ``settle for 50\% of
contract value'' during testing.

The orchestrator exposes both a synchronous \texttt{run()} used by batch scripts
and a \texttt{stream()} generator used by the API. One empirical detail is worth
recording because it caused a real bug: the streamed per-node updates are partial,
and \texttt{bidder\_round} does not carry \texttt{round\_number} in its update
while \texttt{ca\_round} does, so any consumer must read that field from merged
state rather than from the raw update.

\section{From documents to scenarios}
\label{sec:extraction}

Real users upload PDFs and Word files, not Python objects. The extraction agent
takes raw text and produces a validated \texttt{DisputeScenario} in one model call
with a repair retry. Its prompt mirrors the Court agent's discipline in both
directions: do not invent a scoring formula the source does not contain, and do
not summarise away one it does. The second half was added after a real case
lost its published 60/40 weighting during extraction, which silently downgraded
every downstream Court assessment from exact verification to qualitative review.

Two guards were added after failures found by running the system on documents I
had not chosen. A minimum source-length check raises before any model call when
extraction yields under 200 characters, because a scanned PDF with no text layer
previously produced not an error but an entire fictional dispute, complete with a
contract value and a scoring formula. Separately, an unstated contract value must
serialise as \texttt{null} and never as zero, since zero reads downstream as a
genuine \pounds0 contract rather than as missing information.

\section{Structured output and compliance instrumentation}
\label{sec:compliance-method}

Every agent returns JSON validated against a Pydantic schema. Before this project
had instrumentation, a run in which every response needed repair looked identical
to one in which none did, so a counter layer was added around parsing. It records
brace-extraction fallbacks when \texttt{json.loads} fails, field coercions during
validation, repetition regenerations, and outright parse failures. The headline
metric is \emph{structural compliance}: the proportion of structured responses
needing no repair at all. It is scoped per response rather than summed over
events, because one response can trigger several coercions and naively summing
them produces a rate that exceeds one and then goes negative. Repetition retries
are counted but do not mark a response dirty, since repetitiveness is a content
problem rather than a structural one.

\section{Method selection, and what was rejected}
\label{sec:alternatives}

Table~\ref{tab:alternatives} records the significant choices against the
alternatives considered. Three deserve comment.

\begin{table}[htbp]
\centering
\footnotesize
\caption[Design alternatives]{Design decisions against the alternatives considered.
Rows marked \emph{plan change} diverge from the original project plan.}
\label{tab:alternatives}
\begin{tabular}{@{}p{2.6cm}p{3.0cm}p{6.4cm}@{}}
\toprule
\textbf{Decision} & \textbf{Alternative} & \textbf{Reasoning} \\
\midrule
Prompt engineering + context injection & Fine-tuning a domain model &
  No labelled corpus of procurement dispute negotiations exists to fine-tune on;
  building one was not feasible in the time budget. Prompts also stay auditable,
  which matters when the finding under test is fabrication. \\
Llama 3.1 8B \cite{grattafiori2024llama}, self-hosted \emph{(plan change)} & GPT-4o via API &
  Cost, and the fact that client documents are commercially sensitive. Also a
  stronger test: a discipline that holds on an 8B model is not relying on
  capability to mask a weak prompt. \\
Single-pass extraction & Retrieval-augmented generation &
  RAG addresses retrieving the right authority; the failure mode here is inventing
  facts about the case in front of you, which retrieval does not fix. Scoped out
  explicitly rather than left implied. \\
LangGraph state machine & Free-form agent conversation &
  Deterministic turn order and explicit conditional edges make the ablation in
  \S\ref{sec:baselines} possible: you can delete a node and measure it. \\
React frontend \emph{(plan change)} & Streamlit &
  Needed to be shown to a client board, and needed live streaming of an
  in-progress negotiation. \\
JSON logs \emph{(plan change)} & PostgreSQL &
  Every analysis in Chapter~\ref{ch:evaluation} is a batch read over completed
  runs; a database would have added operational surface for no analytical gain. \\
Real judgments only \emph{(plan change)} & Synthetic scenarios &
  The two synthetic scenarios were AI-authored and were deleted, along with their
  evaluation batches, once a real corpus existed. Synthetic disputes cannot
  falsify anything. \\
\bottomrule
\end{tabular}
\end{table}

The move from synthetic to real scenarios is the most consequential. The project
plan proposed synthetic dispute scenarios with real cases as optional validation.
In practice the synthetic pair was written with model assistance, which made it a
circular test: a scenario generated by a language model, assessed by a language
model, has no independent purchase on whether either is right. They were deleted
outright, together with seven batch directories of results derived from them,
rather than retained with a caveat, because the results were still being served
into the frontend's analytics page where no caveat could reach a viewer.

The plan also proposed ablations over temperature and BATNA values. Those were
dropped in favour of the three baselines in Section~\ref{sec:baselines}, on the
judgement that ``does this architecture beat a single prompt'' is the question an
examiner asks first, and it is the question a temperature sweep cannot answer.

\section{Delivery and the industrial deliverable}
\label{sec:risk-screen}

The negotiation simulator is exposed through a FastAPI backend and a
Fusion21-branded React client, with live negotiations streamed over server-sent
events. Inference runs against local Ollama by default and against a tunnelled
University GPU instance when one is reachable, detected at startup and
re-establishable from a button in the interface. That last piece exists because
the tunnel is fragile in practice, and a demonstration that requires a terminal is
not a demonstration.

The component Fusion21 actually asked for is different. Their framing was
prevention: every other part of this system operates on a dispute that has already
been raised. The pre-award challenge risk screen is a rule-based function with no
model call, taking the same 16 profile fields already encoded from Doc 1 and
returning a risk band, a normalised score, and a list of flags. Each flag names the
field that triggered it, the Doc 1 category it traces to, a severity, a confidence
tag, a rationale close-paraphrased from a specific Doc 1 sentence, and a concrete
pre-award mitigation.

Three boundaries are deliberate. It predicts whether a challenge is likely to be
\emph{raised}, not whether the authority would lose one, because the second
question is the Court agent's job and conflating them would smuggle a merits
judgment in through inputs that speak to a bidder's incentive. Every flag carries
a confidence tag, because an authority knows its own documentation quality but is
guessing at a losing bidder's legal representation. And the score is an ordering
device for triage, not a probability: Doc 1 gives directional relationships, not
weights, so the severity weighting is a transparent first pass that the module's
own output says is not calibrated.

\section{Verification discipline as a method}
\label{sec:verification-discipline}

One methodological commitment runs through everything above and is worth stating
as a method rather than as a habit. No number, citation or case fact enters this
project without being checked against a primary source, and that rule was applied
to my own write-ups as well as to model output. It caught real errors in both
directions. An earlier draft of the evaluation asserted that the agents had
negotiated over a case's 68\%/84\% score gap; they had not, because the cached
scenario never contained it, and the sentence had been checked against the raw
case text rather than against what the pipeline actually fed the agents. A
contract value of roughly \pounds1.5 billion attributed to one case turned out, on
verification, to be \pounds112.1 million for the specific procurement at issue.
Both are the same class of error the Court agent is hardened against, occurring in
the human layer, and both are reported here rather than quietly corrected.
